%% Question bank chapter file - generated by bank_tools.py build (do not edit while the bank is still being built:
%% it is regenerated from the per-test data in Question Bank\data). Every question is one bankq block:
%%   \begin{bankq}[id=..., source={...}, paper=.., qno=.., type=..., subtopic={...}, answer={...}, answer_src=key|solved,
%%                 working={...}, difficulty=1-10, diff_parts={K S I C F}, diff_note={...},
%%                 group=<passage id or empty>, also_in={duplicate copies in other tests}]
%% difficulty = min(10, K+S+I+C+F-1): K concepts 1-3, S steps 1-4, I insight 0-2, C computation 0-2, F format 0-1
%%   ...question LaTeX, options, figures (\bfig / \bfigi -> img/)...
%%   \end{bankq}
%% Passages shared by several questions are bankpassage blocks; their questions carry group=<passage id>.
\documentclass[10pt]{article}
\usepackage[a4paper,margin=1.8cm]{geometry}
\usepackage[T1]{fontenc}
\usepackage{amsmath,amssymb,mathtools}
\usepackage{graphicx}
\usepackage[export]{adjustbox}
\usepackage{xcolor,needspace,keyval}
\usepackage{enumitem}
\usepackage{booktabs,array,tabularx}
\usepackage[hidelinks]{hyperref}
\setlength{\parindent}{0pt}
\setlength{\parskip}{3pt}
\makeatletter
\define@key{bq}{id}{\def\bq@id{#1}}\define@key{bq}{source}{\def\bq@src{#1}}\define@key{bq}{paper}{\def\bq@paper{#1}}
\define@key{bq}{qno}{\def\bq@qno{#1}}\define@key{bq}{type}{\def\bq@type{#1}}\define@key{bq}{subtopic}{\def\bq@sub{#1}}
\define@key{bq}{answer}{\def\bq@ans{#1}}\define@key{bq}{answer_src}{\def\bq@asrc{#1}}\define@key{bq}{working}{\def\bq@work{#1}}
\define@key{bq}{group}{\def\bq@group{#1}}\define@key{bq}{also_in}{\def\bq@also{#1}}
\define@key{bq}{difficulty}{\def\bq@diff{#1}}\define@key{bq}{diff_parts}{\def\bq@dparts{#1}}\define@key{bq}{diff_note}{\def\bq@dnote{#1}}
\newenvironment{bankq}[1][]{\def\bq@work{}\def\bq@also{}\def\bq@group{}\def\bq@diff{?}\def\bq@dparts{}\def\bq@dnote{}\setkeys{bq}{#1}%
  \par\needspace{4\baselineskip}\medskip\noindent{\small\textbf{\bq@id}\hfill\color{gray}[\bq@src, Paper \bq@paper, Q\bq@qno]
  \bq@type\ $\cdot$ \bq@sub\ $\cdot$ \textbf{D\bq@diff}}\par\nopagebreak\relax}
  {\par{\small\color{teal}Answer: \textbf{\bq@ans}\ifx\bq@work\@empty\else\ (solved: \bq@work)\fi
  \ifx\bq@also\@empty\else\par\color{gray}Also in: \bq@also\fi}\par}
\define@key{bp}{id}{\def\bp@id{#1}}\define@key{bp}{source}{\def\bp@src{#1}}\define@key{bp}{questions}{\def\bp@qs{#1}}
\newenvironment{bankpassage}[1][]{\setkeys{bp}{#1}\par\needspace{4\baselineskip}\bigskip\noindent
  {\small\textbf{Passage \bp@id}\hfill\color{gray}[\bp@src; questions \bp@qs]}\par\nopagebreak\relax\itshape}{\par}
\makeatother
\newcommand{\bfig}[2]{\par\smallskip{\centering\includegraphics[width=#1pt,max width=0.62\linewidth,max height=0.32\textheight]{img/#2}\par}\smallskip}
\newcommand{\bfigi}[2]{\includegraphics[width=#1pt,max width=\linewidth,max height=0.32\textheight,valign=c]{img/#2}}
\newcommand{\bmissing}[1]{\textit{[figure missing in the source conversion: \detokenize{#1}]}}
\newenvironment{qoptions}{\begin{itemize}[leftmargin=1.8em,itemsep=1pt,topsep=2pt,label={}]}{\end{itemize}}
\newcommand{\opt}[2]{\item[(#1)] #2}
\providecommand{\fwstart}{}\providecommand{\fwend}{}\providecommand{\qfullsplit}{}
\newcolumntype{L}[1]{>{\raggedright\arraybackslash}p{#1}}
\newcolumntype{C}[1]{>{\centering\arraybackslash}p{#1}}
\newcolumntype{R}[1]{>{\raggedleft\arraybackslash}p{#1}}
\providecommand{\sectionhdr}[1]{}\providecommand{\subsectionhdr}[1]{}


\begin{document}
{\Large\bfseries Functions --- Question Bank}\par\smallskip
{\small Mathematics chapter M20. Questions from all converted test papers, grouped by source series. Answers: printed key where it matches the paper, otherwise solved (answer\_src=solved, with working). Exact repeats across papers appear once, with also\_in listing the other copies.}\par\medskip\hrule

\section*{TOAS papers}
\begin{bankq}[id=TOAS01-P1-MATH-S1-Q01, source={TOAS - Test 01 (30 Jan 2022)}, paper=1, qno=1, type=single, subtopic={Range}, answer={C}, answer_src=key, difficulty=3, diff_parts={K1 S2 I1 C0 F0}, diff_note={e^x+e^-x>=2 so two roots need a<-2}]
Let $f(x)=e^x+e^{-x}+a$ has exactly 2 distinct real roots, then
\begin{qoptions}
\opt{A}{$a>0$}
\opt{B}{$a<0$}
\opt{C}{$a<-2$}
\opt{D}{$a\in\phi$}
\end{qoptions}
\end{bankq}
\begin{bankq}[id=TOAS01-P1-MATH-S2-Q08, source={TOAS - Test 01 (30 Jan 2022)}, paper=1, qno=8, type=multi, subtopic={Composition}, answer={A,B,C}, answer_src=key, difficulty=3, diff_parts={K1 S2 I0 C0 F1}, diff_note={Evaluate piecewise f and g at four points}]
Let
\[f(x)=\begin{cases}x^2-3x+4\ ;&x<3\\ x+7\ ;&x\ge3\end{cases}\]
and
\[g(x)=\begin{cases}x+6\ ;&x<4\\ x^2+x+2\ ;&x\ge4\end{cases}\]
then which of the following is/are true -
\begin{qoptions}
\opt{A}{$(f+g)(1)=9$}
\opt{B}{$(f-g)(3.5)=1$}
\opt{C}{$(fg)(0)=24$}
\opt{D}{$\left(\dfrac fg\right)(5)=\dfrac83$}
\end{qoptions}
\end{bankq}
\begin{bankq}[id=TOAS01-P2-MATH-S1-Q16, source={TOAS - Test 01 (30 Jan 2022)}, paper=2, qno=16, type=single, subtopic={Range}, answer={A}, answer_src=key, difficulty=5, diff_parts={K1 S3 I1 C1 F0}, diff_note={Discriminant of y(x^2+4x+2)=11x^2+12x+6 gives range y<=-5 or y>=3}]
Sum of all the integers that do not belong to range of $y=\dfrac{11x^2+12x+6}{x^2+4x+2}$ is -
\begin{qoptions}
\opt{A}{$-7$}
\opt{B}{$-9$}
\opt{C}{$-4$}
\opt{D}{$-2$}
\end{qoptions}
\end{bankq}
\begin{bankq}[id=TOAS01-P2-MATH-S1-Q17, source={TOAS - Test 01 (30 Jan 2022)}, paper=2, qno=17, type=single, subtopic={Composition}, answer={C}, answer_src=key, difficulty=5, diff_parts={K1 S4 I1 C0 F0}, diff_note={Domain and simplification of four function pairs}]
Which of the following pair of function have same graph -
\begin{qoptions}
\opt{A}{$f(x)=\mathrm{sgn}\left(|x|+\dfrac1{|x|}\right)$; $g(x)=\left(\dfrac{x^2+1}{2|x|}\right)$}
\opt{B}{\raggedright $f(x)=\dfrac{(\sin x+\cos x)^2-1}{2\cos^2x}$; $g(x)=\dfrac{2\sin^2x}{1-(\sin x-\cos x)^2}$\par}
\opt{C}{$f(x)=e^{\ln(\mathrm{sgn}\,x^2)}$, $g(x)=\dfrac1{|\mathrm{sgn}(x)|}$}
\opt{D}{$f(x)=\sec^2x-\tan^2x$, $g(x)=1$\\
where [.] represents greatest integer function and sgn (.) represent singum function of $x$}
\end{qoptions}
\end{bankq}
\begin{bankq}[id=TOAS01-P2-MATH-S1-Q18, source={TOAS - Test 01 (30 Jan 2022)}, paper=2, qno=18, type=single, subtopic={Domain}, answer={A}, answer_src=key, difficulty=2, diff_parts={K1 S2 I0 C0 F0}, diff_note={sin x>0 and |x|<=2}]
Domain of function $f(x)=\log_2\sin(x)+\sqrt{4-x^2}$ is -
\begin{qoptions}
\opt{A}{(0, 2]}
\opt{B}{(1, 3]}
\opt{C}{(1, $\pi$]}
\opt{D}{[$-$2, 2]}
\end{qoptions}
\end{bankq}
\begin{bankq}[id=TOAS01-P2-MATH-S1-Q19, source={TOAS - Test 01 (30 Jan 2022)}, paper=2, qno=19, type=single, subtopic={Domain}, answer={C}, answer_src=key, difficulty=6, diff_parts={K2 S3 I1 C1 F0}, diff_note={Product of linear forms in sin, cos equals 3 sin 2x - 4 cos 2x over 2; domain and range}]
Consider $f(x)=\dfrac1{(\sin x+2\cos x)(4\sin x-2\cos x)}$, then -
\begin{qoptions}
\opt{A}{range of $f(x)$ is $\left[-\dfrac1{\sqrt5},\ \dfrac1{\sqrt5}\right]$}
\opt{B}{domain of $f(x)$ is $R-(n\pi+\alpha)\ \forall\ \alpha=\tan^{-1}\left(\dfrac43\right)$, $n\in I$}
\opt{C}{domain of $f(x)$ is $R-\left(\dfrac{n\pi}2+\dfrac\alpha2\right)\ \forall\ \alpha=\tan^{-1}\dfrac43$, $n\in I$}
\opt{D}{range of $f(x)$ is $\left(-\infty,\ -\dfrac1{\sqrt5}\right]\cup\left[\dfrac1{\sqrt5},\ \infty\right)$}
\end{qoptions}
\end{bankq}
\begin{bankq}[id=TOAS01-P2-MATH-S1-Q20, source={TOAS - Test 01 (30 Jan 2022)}, paper=2, qno=20, type=single, subtopic={Modulus and greatest integer}, answer={A}, answer_src=key, difficulty=1, diff_parts={K1 S1 I0 C0 F0}, diff_note={sqrt(sgn|x|) is 1 for x not 0 and 0 at 0}]
Which of the following best suits the graph of $f(x)=\sqrt{\mathrm{Sgn}|x|}$
\begin{qoptions}
\opt{A}{\bfigi{92.0}{TOAS01_p2_math_s1_q20_A.png}}
\opt{B}{\bfigi{92.8}{TOAS01_p2_math_s1_q20_B.png}}
\opt{C}{\bfigi{96.0}{TOAS01_p2_math_s1_q20_C.png}}
\opt{D}{\bfigi{93.0}{TOAS01_p2_math_s1_q20_D.png}}
\end{qoptions}
\end{bankq}
\begin{bankq}[id=TOAS02-P1-MATH-S1-Q02, source={TOAS - Test 02 (13 Feb 2022)}, paper=1, qno=2, type=single, subtopic={Range}, answer={B}, answer_src=key, difficulty=5, diff_parts={K2 S3 I1 C0 F0}, diff_note={Range from the discriminant forces a=0 and b=-4}]
Consider $f:R\to R$, $f(x)=\dfrac{3x^2+ax+b}{x^2+1}$ where $a,b\in R$. If the range of this function is $[-4,3)$ then $a+b$ is equal to
\begin{qoptions}
\opt{A}{0}
\opt{B}{$-4$}
\opt{C}{$-1$}
\opt{D}{3}
\end{qoptions}
\end{bankq}
\begin{bankq}[id=TOAS02-P1-MATH-S1-Q03, source={TOAS - Test 02 (13 Feb 2022)}, paper=1, qno=3, type=single, subtopic={Composition}, answer={C}, answer_src=key, difficulty=5, diff_parts={K1 S3 I2 C0 F0}, diff_note={Chain of inequalities forces f(x+1)=f(x)+1, then g(2021)}]
Given that $f(x)$ is a function defined on $\mathbf R$, satisfying $f(1)=1$ and for any $x\in\mathbf R$, $f(x+5)\ge f(x)+5$ and $f(x+1)\le f(x)+1$. If $g(x)=f(x)+5-x$, then $g(2021)$ is equal to -
\begin{qoptions}
\opt{A}{3}
\opt{B}{4}
\opt{C}{5}
\opt{D}{6}
\end{qoptions}
\end{bankq}
\begin{bankq}[id=TOAS02-P1-MATH-S1-Q04, source={TOAS - Test 02 (13 Feb 2022)}, paper=1, qno=4, type=single, subtopic={Range}, answer={C}, answer_src=key, difficulty=5, diff_parts={K2 S3 I1 C0 F0}, diff_note={Functional equation gives f(x)=x^2+x+1, test four statements}]
Let $f:R\to R$ and $f(x)$ be a polynomial with real coefficient such that $(x^2+x+1)\,f(x-1)=(x^2-x+1)\,f(x)\ \forall\ x\in R$ and $f(0)=1$, then which of the following is false -
\begin{qoptions}
\opt{A}{$f(x)$ is many one function}
\opt{B}{$f(x)$ has no real roots}
\opt{C}{range of $f(x)$ is $\left[\dfrac12,\infty\right)$}
\opt{D}{$f(x)$ is neither even nor odd}
\end{qoptions}
\end{bankq}
\begin{bankq}[id=TOAS02-P1-MATH-S1-Q05, source={TOAS - Test 02 (13 Feb 2022)}, paper=1, qno=5, type=single, subtopic={Modulus and greatest integer}, answer={C}, answer_src=key, difficulty=5, diff_parts={K2 S3 I1 C0 F0}, diff_note={a^3+b^3+c^3-3abc factorisation with integer parts, minimum sum of floors}]
If $[x]^3+[y]^3+[z]^3-3[x][y][z]=4$; where $x,y,z>0$ (where [.] denotes GIF); then minimum possible value of $(x+y+z)$ is
\begin{qoptions}
\opt{A}{2}
\opt{B}{3}
\opt{C}{4}
\opt{D}{6}
\end{qoptions}
\end{bankq}
\begin{bankq}[id=TOAS02-P1-MATH-S2-Q12, source={TOAS - Test 02 (13 Feb 2022)}, paper=1, qno=12, type=multi, subtopic={Modulus and greatest integer}, answer={A,B}, answer_src=key, difficulty=7, diff_parts={K2 S4 I1 C0 F1}, diff_note={Graph transformations with sgn and modulus for four expressions}]
The graph of the function $y=f(x)$ is shown in the figure. Then which one of the following graph(s) of corresponding curve is/are correct ?
\bfig{198.5}{TOAS02_p1_math_s2_q12_1.png}
\begin{qoptions}
\opt{A}{$|y|=\mathrm{sgn}(f(x))$\\ \bfigi{103.1}{TOAS02_p1_math_s2_q12_2.png}}
\opt{B}{$|y|=\mathrm{sgn}(-f(|x|))$\\ \bfigi{132.3}{TOAS02_p1_math_s2_q12_3.png}}
\opt{C}{$|y|=|f(|x|)|$\\ \bfigi{139.9}{TOAS02_p1_math_s2_q12_4.png}}
\opt{D}{$y=x\,\mathrm{sgn}(f(x))$\\ \bfigi{183.0}{TOAS02_p1_math_s2_q12_5.png}}
\end{qoptions}
\end{bankq}
\begin{bankq}[id=TOAS02-P1-MATH-S3-Q04, source={TOAS - Test 02 (13 Feb 2022)}, paper=1, qno=4, type=numerical, subtopic={Periodicity}, answer={256}, answer_src=solved, working={Periods are $2\pi\cdot2^k$ for sine and $\pi\cdot2^k$ for tangent terms; largest ($\sin\frac x{2^7}$, $\tan\frac x{2^8}$) is $256\pi$ and all others divide it}, difficulty=5, diff_parts={K1 S3 I1 C0 F1}, diff_note={LCM of periods of sine and tangent terms: 256 pi}]
If the fundamental period of the function $f(x)=\sin x+\tan\dfrac x2+\sin\dfrac x{2^2}+\tan\dfrac x{2^3}+\dots+\sin\dfrac x{2^7}+\tan\dfrac x{2^8}$ is $k\pi$, then $k$ is equal to -
\end{bankq}
\begin{bankq}[id=TOAS03-P1-MATH-S1-Q03, source={TOAS - Test 03 (27 Feb 2022)}, paper=1, qno=3, type=single, subtopic={Range}, answer={A}, answer_src=key, difficulty=6, diff_parts={K2 S3 I2 C0 F0}, diff_note={Functional equation solved by f(x)=e^x-1 via f(x)=f(x/2^n) product limit}]
Let $f:\mathbb R\to\mathbb R$ be a non-constant continuous function such that $f(2x)=(e^x+1)f(x)$, $f(0)=0$ \& $f'(0)=1$, then range of $f(x)$ is-
\begin{qoptions}
\opt{A}{$(-1,\infty)$}
\opt{B}{$[0,\infty)$}
\opt{C}{$[1,\infty)$}
\opt{D}{$(1,\infty)$}
\end{qoptions}
\end{bankq}
\begin{bankq}[id=TOAS03-P1-MATH-S2-Q09, source={TOAS - Test 03 (27 Feb 2022)}, paper=1, qno=9, type=multi, subtopic={Domain}, answer={A,C,D}, answer_src=key, difficulty=5, diff_parts={K1 S3 I1 C0 F1}, diff_note={Solve each radicand in |x| and intersect to get the domain}]
If $f(x)=\sqrt{|x|^2-5|x|+6}+\sqrt{8+2|x|-|x|^2}$, then $x$ can be-
\begin{qoptions}
\opt{A}{$[-4,-3]$}
\opt{B}{$[-3,-2]$}
\opt{C}{$[-2,2]$}
\opt{D}{$[3,4]$}
\end{qoptions}
\end{bankq}
\begin{bankq}[id=TOAS03-P1-MATH-S3-Q01, source={TOAS - Test 03 (27 Feb 2022)}, paper=1, qno=1, type=numerical, subtopic={Composition}, answer={2.50}, answer_src=key, difficulty=5, diff_parts={K1 S3 I1 C0 F1}, diff_note={Involution g(g(x))=x gives f(x)=2g(x)-x}]
Let $f$ be a function satisfying the functional equation $f(x)+2f\left(\dfrac{2x+1}{x-2}\right)=3x,\ x\neq2$, then the value of $\dfrac12\left[\left|\dfrac{f(3)}{2f(7)}\right|\right]$ is [Note : [.] denotes greatest integer function]
\end{bankq}
\begin{bankq}[id=TOAS03-P1-MATH-S3-Q02, source={TOAS - Test 03 (27 Feb 2022)}, paper=1, qno=2, type=numerical, subtopic={Domain}, answer={4.00}, answer_src=key, difficulty=5, diff_parts={K1 S3 I1 C0 F1}, diff_note={Domain integers from radicand and denominator conditions}]
The number of integers in the domain of $f(x)=\dfrac{\sqrt{(x^2+kx+k)}}{x^2+k}$ is
\end{bankq}
\begin{bankq}[id=TOAS03-P2-MATH-S1-Q01, source={TOAS - Test 03 (27 Feb 2022)}, paper=2, qno=1, type=multi, subtopic={Composition}, answer={C,D}, answer_src=key, difficulty=5, diff_parts={K1 S3 I1 C0 F1}, diff_note={g(x) simplifies to the constant 2, domain not symmetric so neither even nor odd, bounded}]
Let $f(x)=\dfrac{2x+1}{x},\ x\neq0$, where \[g(x)=f(f(x))-\dfrac1{f(x)},\quad x\neq0,-\dfrac12,\] then identify the correct statement(s) -
\begin{qoptions}
\opt{A}{$g(x)$ is an odd function}
\opt{B}{$g(x)$ is an even function}
\opt{C}{$g(x)$ is neither even nor odd}
\opt{D}{$g(x)$ is a bounded function}
\end{qoptions}
\end{bankq}
\begin{bankq}[id=TOAS03-P2-MATH-S1-Q05, source={TOAS - Test 03 (27 Feb 2022)}, paper=2, qno=5, type=multi, subtopic={Range}, answer={A,B,C,D}, answer_src=key, difficulty=4, diff_parts={K1 S3 I0 C0 F1}, diff_note={f takes values 2, 0, -2 on three intervals}]
Consider a function $f(x)=\dfrac{|x+1|}{x+1}+\dfrac{x-1}{|x-1|}$, then choose correct option(s) regarding function $f$
\begin{qoptions}
\opt{A}{many one}
\opt{B}{odd function}
\opt{C}{bounded function}
\opt{D}{range contain exactly three integers}
\end{qoptions}
\end{bankq}
\begin{bankq}[id=TOAS03-P2-MATH-S2-Q01, source={TOAS - Test 03 (27 Feb 2022)}, paper=2, qno=1, type=numerical, subtopic={Inverse}, answer={0.50}, answer_src=key, difficulty=5, diff_parts={K1 S3 I1 C0 F1}, diff_note={f is an involution so f^-1=f; equation reduces to 2x^2-8x+8=0}]
Let $f:\mathbb R-\left\{\dfrac12\right\}\to\mathbb R-\left\{\dfrac12\right\}$, $f(x)=\dfrac{x-2}{2x-1}$ be a function such that $x=m$ is the solution of \[f(x)+2f^{-1}(x)+2=f(f(x)),\] then $\dfrac m4$ is equal to
\end{bankq}
\begin{bankq}[id=TOAS03-P2-MATH-S2-Q02, source={TOAS - Test 03 (27 Feb 2022)}, paper=2, qno=2, type=numerical, subtopic={Domain}, answer={0.00}, answer_src=key, difficulty=3, diff_parts={K1 S2 I0 C0 F1}, diff_note={ln of inverse trig needs positive arguments: x in (0,1)}]
The number of integers in the domain of the function $f(x)=e^{\ell n(\sin^{-1}x)}+e^{\ell n(\cos^{-1}x)}$ is
\end{bankq}
\begin{bankq}[id=TOAS03-P2-MATH-S3-Q01, source={TOAS - Test 03 (27 Feb 2022)}, paper=2, qno=1, type=numerical, subtopic={Periodicity}, answer={8}, answer_src=key, difficulty=6, diff_parts={K2 S3 I1 C0 F1}, diff_note={f(x+3)=-f(x) gives period 6; even function relates f(-5), f(1) and f(-13)}]
Let $f(x)$ be an even function satisfying $f(x)+f(x+2)=f(x+1)\ \forall\ x\in\mathbb R$, where $f(-13)=4$, then $f(-5)+f(1)$ is equal to
\end{bankq}
\begin{bankq}[id=TOAS03-P2-MATH-S3-Q02, source={TOAS - Test 03 (27 Feb 2022)}, paper=2, qno=2, type=numerical, subtopic={Inverse}, answer={1}, answer_src=key, difficulty=4, diff_parts={K1 S2 I1 C0 F1}, diff_note={A Mobius map is an involution when the trace is zero: b=-1}]
If $f(x)=\dfrac{x+a}{x+b}$ and $f:\mathbb R-\{-b\}\to\mathbb R-\{1\}$ is such that $f(x)=f^{-1}(x)$, then value of $|b|$ must be
\end{bankq}
\begin{bankq}[id=TOAS04-P1-MATH-S2-Q07, source={TOAS - Test 04 (15 Mar 2022)}, paper=1, qno=7, type=multi, subtopic={Composition}, answer={A,B,D}, answer_src=key, difficulty=6, diff_parts={K1 S3 I2 C0 F1}, diff_note={Bijective composition forces h injective, f surjective; g need not be onto}]
$f:\mathbb R\to\mathbb R$, $g:\mathbb R\to\mathbb R$, $h:\mathbb R\to\mathbb R$, defined such that $f\circ g\circ h:\mathbb R\to\mathbb R$ is a bijective function then :
\begin{qoptions}
\opt{A}{$h$ must be injective}
\opt{B}{$g$ may be many one into}
\opt{C}{$g$ must be surjective}
\opt{D}{$f$ must be surjective}
\end{qoptions}
\end{bankq}
\begin{bankq}[id=TOAS04-P1-MATH-S2-Q09, source={TOAS - Test 04 (15 Mar 2022)}, paper=1, qno=9, type=multi, subtopic={Periodicity}, answer={A,B,C,D}, answer_src=key, difficulty=6, diff_parts={K2 S3 I1 C0 F1}, diff_note={Counterexamples to four claims on periods, inverse and additive functions}]
Which of the following are incorrect ?
\begin{qoptions}
\opt{A}{Let $f$ be non constant periodic function with periods (need not be fundamental period) $T_1$ as well as $T_2$ then $\dfrac{T_1}{T_2}$ must be rational}
\opt{B}{Let $f$ and $g$ be periodic with periods $T_1$ and $T_2$ respectively such that LCM $(T_1,T_2)$ does not exist then $f+g$ must be non periodic}
\opt{C}{Let $f:\mathbb R\to\mathbb R$ be an invertible function then $f(x)=x\Leftrightarrow f(x)=f^{-1}(x)$}
\opt{D}{Let $f:\mathbb R\to\mathbb R$ be a function such that $f(x+y)=f(x)+f(y)\ \forall\ x,y\in\mathbb R$ then $f(x)$ must be $kx$ for some real number $k\in\mathbb R$.}
\end{qoptions}
\end{bankq}
\begin{bankq}[id=TOAS04-P2-MATH-S1-Q06, source={TOAS - Test 04 (15 Mar 2022)}, paper=2, qno=6, type=multi, subtopic={Range}, answer={A,B,C}, answer_src=key, difficulty=6, diff_parts={K1 S3 I2 C0 F1}, diff_note={Telescoping bound forces f constant}]
If a function $f:R\to R$ satisfies $\left|\displaystyle\sum_{k=1}^{n}3^k\left(f(x+ky)-f(x-ky)\right)\right|\le1$ for every $n\in N$ and for all $x,y\in R$, then which of the following is/are true
\begin{qoptions}
\opt{A}{$\left|3^n\left(f(x+ny)-f(x-ny)\right)\right|\le2$}
\opt{B}{$|f(u)-f(v)|\le\dfrac{2}{3^n}$}
\opt{C}{$f(x)$ is a constant function}
\opt{D}{$f(x)$ is an odd function}
\end{qoptions}
\end{bankq}
\begin{bankpassage}[id=TOAS04-P2-MATH-S3-G5-6, source={TOAS - Test 04 (15 Mar 2022)}, questions={5-6}]
Consider the function $f:A\to[0,5]$ $f(x)=\sqrt{9-(x-1)(x-3)(x-5)(x-7)}$.\\
(Where set $A$ represents all values of $x$ for which $f(x)$ is defined.)
\end{bankpassage}
\begin{bankq}[id=TOAS04-P2-MATH-S3-Q05, source={TOAS - Test 04 (15 Mar 2022)}, paper=2, qno=5, type=numerical, subtopic={Domain}, answer={7.00}, answer_src=key, difficulty=5, diff_parts={K1 S3 I1 C0 F1}, diff_note={u=(x-4)^2 reduces the radicand to u(u-10)<=0: integers 1 to 7}, group=TOAS04-P2-MATH-S3-G5-6]
In the domain of the function, number of possible integral values of $x$ is equal to
\end{bankq}
\begin{bankq}[id=TOAS04-P2-MATH-S3-Q06, source={TOAS - Test 04 (15 Mar 2022)}, paper=2, qno=6, type=numerical, subtopic={Domain}, answer={2.00}, answer_src=key, difficulty=5, diff_parts={K1 S3 I1 C0 F1}, diff_note={Function is many-one and onto [0,5]}, group=TOAS04-P2-MATH-S3-G5-6]
Among the following, the number of correct statements is\\
(i) function is many one \quad (ii) function is one-one\\
(iii) function is onto \quad (iv) function is into
\end{bankq}
\begin{bankq}[id=TOAS04-P2-MATH-S4-Q07, source={TOAS - Test 04 (15 Mar 2022)}, paper=2, qno=7, type=numerical, subtopic={Composition}, answer={561}, answer_src=key, difficulty=6, diff_parts={K1 S3 I1 C1 F1}, diff_note={f(x)-f(x-2)=2x-1 summed to reach f(93) then f(94)}]
Let $f(x)$ be a function satisfying $f(x)+f(x-1)=x^2$, $\forall\,x\in R$. If $f(19)=94$ then the remainder when $f(94)$ is divided by 1000 is
\end{bankq}
\begin{bankq}[id=TOAS04-P2-MATH-S4-Q09, source={TOAS - Test 04 (15 Mar 2022)}, paper=2, qno=9, type=numerical, subtopic={Composition}, answer={429}, answer_src=key, difficulty=7, diff_parts={K1 S3 I2 C1 F1}, diff_note={f(3x)=3f(x): reduce 2001 to [1,3] and locate the smallest preimage of 186}, also_in={TOAS - Test 09 (30 May 2022) P1 Q7}]
A function $f$ satisfies $f(3x)=3f(x)$ $\forall\,x\in[1,\infty)$ and $f(x)=1-|x-2|$ for $x\in[1,3]$. The smallest $x$ for which $f(x)=f(2001)$ is
\end{bankq}
\begin{bankq}[id=TOAS05-P1-MATH-S1-Q01, source={TOAS - Test 05 (03 Apr 2022)}, paper=1, qno=1, type=multi, subtopic={Composition}, answer={A,B,C,D}, answer_src=key, difficulty=5, diff_parts={K1 S2 I2 C0 F1}, diff_note={Functional equation forces f(x)f(y)=0, so f is identically zero}]
Consider $f:R\to R$ satisfying the equation $f(x+y)=f(x-y)+f(x).f(y)$ $\forall\ x,y\in R$.\
Identify the correct statement(s)
\begin{qoptions}
\opt{A}{$f(2014)+f(2015)=f(2016)$}
\opt{B}{$\sin^{-1}(f(2))+\cos^{-1}(f(3))=\dfrac\pi2$}
\opt{C}{$f(2014)-f(2015)=f(2016)$}
\opt{D}{$\sin^{-1}(f(3))+\cos^{-1}(f(2))=\dfrac\pi2$}
\end{qoptions}
\end{bankq}
\begin{bankq}[id=TOAS05-P1-MATH-S1-Q05, source={TOAS - Test 05 (03 Apr 2022)}, paper=1, qno=5, type=multi, subtopic={Composition}, answer={B,C,D}, answer_src=key, difficulty=6, diff_parts={K1 S4 I1 C0 F1}, diff_note={Even/odd parity of four function descriptions}]
Which of the following are even functions in their domain
\begin{qoptions}
\opt{A}{$y(x)=\dfrac{(2\sin x-1)(\sin^2x-1)}{(2\sin x-1)}$}
\opt{B}{$y(x)=f\mathrm{ogoh}(x)$ (where $f$ is neither even nor odd, $g$ is even and $h$ is an odd function, all functions are defined for all $x\in R$)}
\opt{C}{$y(x)$ is defined for all real $x$ and $y(x)=|x-1|$ for $x\in[0,2]$, and $y(x)$ is periodic with period 2.}
\opt{D}{$y(x)=\dfrac{1+e^{x-1}\left(\cos(1-x)+\cos(1+x)\right)+e^{2x}}{e^x}$}
\end{qoptions}
\end{bankq}
\begin{bankq}[id=TOAS05-P2-MATH-S1-Q01, source={TOAS - Test 05 (03 Apr 2022)}, paper=2, qno=1, type=single, subtopic={Composition}, answer={C}, answer_src=key, difficulty=5, diff_parts={K1 S3 I1 C1 F0}, diff_note={Substitute x -> 1/x in the functional relation and solve for f(x+1)}]
Let $f$ be a function defined for all $x>1$ such that $4f(x^{-1}+1)+8f(x+1)=\log_{12}x$. Then $4(f(10)+f(13)+f(17))$ is-
\begin{qoptions}
\opt{A}{1}
\opt{B}{2}
\opt{C}{3}
\opt{D}{4}
\end{qoptions}
\end{bankq}
\begin{bankq}[id=TOAS05-P2-MATH-S1-Q03, source={TOAS - Test 05 (03 Apr 2022)}, paper=2, qno=3, type=single, subtopic={Inverse}, answer={B}, answer_src=key, difficulty=3, diff_parts={K1 S2 I1 C0 F0}, diff_note={A Mobius map is its own inverse when trace is zero: b=-1}]
If the function $f:R-\{-b\}\to R-\{1\}$ is defined by $f(x)=\dfrac{(x+a)}{(x+b)}$ ($a\ne b$) is inverse of itself, then $b+2$ is
\begin{qoptions}
\opt{A}{2}
\opt{B}{1}
\opt{C}{5}
\opt{D}{8}
\end{qoptions}
\end{bankq}
\begin{bankq}[id=TOAS05-P2-MATH-S1-Q04, source={TOAS - Test 05 (03 Apr 2022)}, paper=2, qno=4, type=single, subtopic={Composition}, answer={D}, answer_src=key, difficulty=4, diff_parts={K1 S2 I2 C0 F0}, diff_note={f(f(x))=x so even number of iterations returns 1}]
If $f(x)=\sqrt[5]{5-x^5}$, then \[\underbrace{f(f(f.\ldots.(1)))}_{2016\ \text{times}\ f}\] is equal to -
\begin{qoptions}
\opt{A}{$4^{1/5}$}
\opt{B}{$4^{\frac{2016}{5}}$}
\opt{C}{$4^{\frac1{10080}}$}
\opt{D}{1}
\end{qoptions}
\end{bankq}
\begin{bankq}[id=TOAS05-P2-MATH-S2-Q05, source={TOAS - Test 05 (03 Apr 2022)}, paper=2, qno=5, type=numerical, subtopic={Domain}, answer={0}, answer_src=solved, working={$3^y>0$ needs $4x^2-1>x^4$, $x^2\in(2-\sqrt3,2+\sqrt3)$: domain symmetric, $(-b_0,-a_0)\cup(a_0,b_0)$, sum of endpoints 0}, difficulty=6, diff_parts={K1 S3 I1 C1 F1}, diff_note={Domain from 4x^2-1 > x^4: two symmetric intervals}]
The whole domain of $y=f(x)$ defined implicitly by the equation $3^y+2^{x^4}=2^{4x^2-1}$ is $(a,b)\cup(c,d)$, then $a+b+c+d$ is
\end{bankq}
\begin{bankq}[id=TOAS06-P1-MATH-S1-Q01, source={TOAS - Test 06 (17 Apr 2022)}, paper=1, qno=1, type=single, subtopic={Modulus and greatest integer}, answer={A}, answer_src=key, difficulty=4, diff_parts={K1 S3 I1 C0 F0}, diff_note={[x/5pi] is 0 for x>=0 and -1 for x<0, so f(x)=2|x|cos x is even}]
If $f:[-2\pi,2\pi]\to\mathbb R$ is given by $f(x)=\dfrac{x\cos x}{\left[\frac{x}{5\pi}\right]+\frac12}$ where $[x]$ denotes greatest integer function then $f(x)$
\begin{qoptions}
\opt{A}{is an even function}
\opt{B}{is an odd function}
\opt{C}{is neither even nor odd}
\opt{D}{satisfies $f(x)-f(-x)=2$}
\end{qoptions}
\end{bankq}
\begin{bankq}[id=TOAS06-P1-MATH-S3-Q06, source={TOAS - Test 06 (17 Apr 2022)}, paper=1, qno=6, type=numerical, subtopic={Composition}, answer={3.00}, answer_src=key, difficulty=5, diff_parts={K1 S3 I1 C0 F1}, diff_note={Cyclic substitution x->1-1/x: solve the three-equation system for f(1/2)}]
If $f:R-\{0,1\}\to R$ be a function satisfying $f(x)+f\left(1-\dfrac1x\right)=\tan^{-1}x$ and $N=f\left(\dfrac12\right)$, then the value of $\left(\dfrac{8N}{\pi}\right)$ is
\end{bankq}
\begin{bankq}[id=TOAS07-P1-MATH-S1-Q03, source={TOAS - Test 07 (01 May 2022)}, paper=1, qno=3, type=single, subtopic={Composition}, answer={B}, answer_src=key, difficulty=8, diff_parts={K2 S4 I2 C1 F0}, diff_note={Substitute x -> -x to solve for f and test four properties}]
Let $f(x)$ be a function satisfying the condition $f(|x-3|)+2f(|x+3|)=x^2+2x\ \forall\ x\in R$, then which of the following is not correct
\begin{qoptions}
\opt{A}{minimum value of $f(x)$ is $-3$}
\opt{B}{$f(x)$ is one-one function}
\opt{C}{equation $|f(x)|=2$ has four distinct solutions}
\opt{D}{$|f(x)|$ is non-differentiable at two points}
\end{qoptions}
\end{bankq}
\begin{bankq}[id=TOAS07-P1-MATH-S1-Q04, source={TOAS - Test 07 (01 May 2022)}, paper=1, qno=4, type=single, subtopic={Inverse}, answer={C}, answer_src=key, difficulty=4, diff_parts={K1 S3 I1 C0 F0}, diff_note={Invert h, g, f in turn at 4}]
For three invertible functions, $f(x)=\dfrac{x-3}{x+3}$, $g(x)=2x$ and $h(x)=2^{\frac{x+3}{x+2}}$, value of $f^{-1}(g^{-1}(h^{-1}(4)))$ is -
\begin{qoptions}
\opt{A}{$-\dfrac12$}
\opt{B}{$\dfrac12$}
\opt{C}{1}
\opt{D}{$-1$}
\end{qoptions}
\end{bankq}
\begin{bankq}[id=TOAS07-P1-MATH-S2-Q06, source={TOAS - Test 07 (01 May 2022)}, paper=1, qno=6, type=multi, subtopic={Inverse}, answer={B,C,D}, answer_src=key, difficulty=6, diff_parts={K2 S3 I1 C0 F1}, diff_note={Inverse of log2(2^(x^3+x)+4): g(3), g'(3), limit and range}]
Let $f:R\to A$ be an invertible function defined by $f(x)=\log_2\left(2^{x^3+x}+4\right)$. If $g(x)$ be its inverse function, then -
\begin{qoptions}
\opt{A}{$g(3)=2$}
\opt{B}{$g'(3)=\dfrac12$}
\opt{C}{$\displaystyle\lim_{x\to3}(g(x))^{\frac{\ln4}{x-3}}=2$}
\opt{D}{set $A$ is $(2,\infty)$}
\end{qoptions}
\end{bankq}
\begin{bankq}[id=TOAS07-P1-MATH-S2-Q07, source={TOAS - Test 07 (01 May 2022)}, paper=1, qno=7, type=multi, subtopic={Inverse}, answer={A,D}, answer_src=key, difficulty=8, diff_parts={K2 S3 I2 C1 F1}, diff_note={Inverse of a monotone function: derivative via inverse rule at specific points}]
If the function $f(x)=-4e^{\frac{1-x}{2}}+1+x+\dfrac{x^2}{2}+\dfrac{x^3}{3}$ and $g(x)=f^{-1}(x)$, then -
\begin{qoptions}
\opt{A}{$g'\left(-\dfrac76\right)=\dfrac15$}
\opt{B}{$g'\left(\dfrac76\right)=\dfrac15$}
\opt{C}{$g(1)=-\dfrac76$}
\opt{D}{$g\left(-\dfrac76\right)=1$}
\end{qoptions}
\end{bankq}
\begin{bankpassage}[id=TOAS07-P1-MATH-S3-G3-4, source={TOAS - Test 07 (01 May 2022)}, questions={3-4}]
Let $f:R\to[0,\infty)$ be a derivable function which satisfies $f(x^2+y^2)=f(2xy)+f(x^2-y^2)\ \forall\ x,y\in R$. Also $f'(1)=2$.
\end{bankpassage}
\begin{bankq}[id=TOAS07-P1-MATH-S3-Q03, source={TOAS - Test 07 (01 May 2022)}, paper=1, qno=3, type=numerical, subtopic={Composition}, answer={25.00}, answer_src=key, difficulty=6, diff_parts={K1 S3 I2 C0 F1}, diff_note={Functional equation f=x^2 form: f(3)+f(4)}, group=TOAS07-P1-MATH-S3-G3-4]
$f(3)+f(4)$ is equal to -
\end{bankq}
\begin{bankq}[id=TOAS07-P1-MATH-S3-Q04, source={TOAS - Test 07 (01 May 2022)}, paper=1, qno=4, type=numerical, subtopic={Composition}, answer={4.00}, answer_src=key, difficulty=5, diff_parts={K1 S3 I1 C0 F1}, diff_note={Solutions of cos^2=sin^2 in [0,2 pi]}, group=TOAS07-P1-MATH-S3-G3-4]
Number of solutions of $f(\cos\theta)=f(\sin\theta)$ in $[0,2\pi]$ -
\end{bankq}
\begin{bankq}[id=TOAS07-P1-MATH-S4-Q07, source={TOAS - Test 07 (01 May 2022)}, paper=1, qno=7, type=numerical, subtopic={Modulus and greatest integer}, answer={4}, answer_src=key, difficulty=6, diff_parts={K2 S3 I1 C0 F1}, diff_note={Floor-function cubic identity gives integer triples; maximise (a+b)^c}]
If $[x]^3+[y]^3+[z]^3-3[x][y][z]=4$ where $x,y,z>0$ then maximum possible value of $([x]+[y])^{[z]}$ is (where $[.]$ represents greatest integer function)
\end{bankq}
\begin{bankq}[id=TOAS07-P2-MATH-S1-Q03, source={TOAS - Test 07 (01 May 2022)}, paper=2, qno=3, type=multi, subtopic={Composition}, answer={A,C,D}, answer_src=key, difficulty=6, diff_parts={K1 S3 I2 C0 F1}, diff_note={Additive-type functional equation gives f linear with f'(x) constant}]
Let $f(x+y)+f(x-y)=2f(x)\ \forall\ x,y\in R$, then which of the following is/are be correct ($f(0)=0$ \& $f$ is differentiable $\forall\ x\in R$)
\begin{qoptions}
\opt{A}{$f$ must be odd}
\opt{B}{$f$ must be even}
\opt{C}{$f'(1)=f'(2)$}
\opt{D}{$f'(2)=f'(3)=f'(4)$}
\end{qoptions}
\end{bankq}
\begin{bankq}[id=TOAS07-P2-MATH-S1-Q06, source={TOAS - Test 07 (01 May 2022)}, paper=2, qno=6, type=multi, subtopic={Range}, answer={B,C}, answer_src=key, difficulty=6, diff_parts={K2 S3 I1 C0 F1}, diff_note={Horizontal line cutting the piecewise graph at four points}]
Consider
\[ f(x)=\begin{cases}x^2-2,&x<0\\x,&0\leqslant x<2\\(x-2)^2,&2\leqslant x<4\\x-4,&x\geqslant4\end{cases} \]
If the equation $f(x)=k$ has four roots, then $k$ can be equal to -
\begin{qoptions}
\opt{A}{$-2$}
\opt{B}{0}
\opt{C}{1}
\opt{D}{2}
\end{qoptions}
\end{bankq}
\begin{bankpassage}[id=TOAS07-P2-MATH-S2-G9-10, source={TOAS - Test 07 (01 May 2022)}, questions={9-10}]
Consider the functions $f(x)=\dfrac{2^x-1}{2^x+1}$; $g(x)=\sin^{-1}x$ , $gof:R\to\left[-\dfrac\pi2,\dfrac\pi2\right]$ where $gof(x)=g(f(x))$.\\
On the basis of above information, answer the following questions :
\end{bankpassage}
\begin{bankq}[id=TOAS07-P2-MATH-S2-Q09, source={TOAS - Test 07 (01 May 2022)}, paper=2, qno=9, type=single, subtopic={Inverse}, answer={B}, answer_src=key, difficulty=3, diff_parts={K1 S2 I1 C0 F0}, diff_note={f is bijective onto (-1,1) so gof is one-one into}, group=TOAS07-P2-MATH-S2-G9-10]
$gof$ is -
\begin{qoptions}
\opt{A}{one-one onto}
\opt{B}{one-one into}
\opt{C}{many one}
\opt{D}{many one into}
\end{qoptions}
\end{bankq}
\begin{bankq}[id=TOAS07-P2-MATH-S2-Q10, source={TOAS - Test 07 (01 May 2022)}, paper=2, qno=10, type=single, subtopic={Inverse}, answer={A}, answer_src=key, difficulty=4, diff_parts={K1 S3 I1 C0 F0}, diff_note={Limit of f(g(f(x))) as x to -infinity and floor}, group=TOAS07-P2-MATH-S2-G9-10]
$\left[\displaystyle\lim_{x\to-\infty}f(g(f(x)))\right]$ is (where $[.]$ greatest integer function) -
\begin{qoptions}
\opt{A}{$-1$}
\opt{B}{$-2$}
\opt{C}{0}
\opt{D}{1}
\end{qoptions}
\end{bankq}
\begin{bankq}[id=TOAS08-P1-MATH-S3-Q03, source={TOAS - Test 08 (15 May 2022)}, paper=1, qno=3, type=numerical, subtopic={Composition}, answer={1.50}, answer_src=solved, working={$y=1$ gives $f(x)=3^x$; the limit is $(3^x+3^{-x}-2)/x^2\to(\ln3)^2$ so $k=3$, $k/2=1.5$}, difficulty=8, diff_parts={K2 S3 I2 C1 F1}, diff_note={Functional equation gives f=3^x, then the limit}]
Let $f(x)$ be a real valued function satisfying $2f(xy)=(f(x))^y+(f(y))^x\ \forall x,y\in R$, $f(1)=3$ and $\displaystyle\lim_{x\to0}\dfrac{f(x)+f(-x)-f(f(0))+1}{x^2}$ is equal to $(\ln k)^2$ where $k$ is natural number, then the value of $\dfrac k2$ is
\end{bankq}
\begin{bankq}[id=TOAS08-P1-MATH-S3-Q08, source={TOAS - Test 08 (15 May 2022)}, paper=1, qno=8, type=numerical, subtopic={Composition}, answer={0.75}, answer_src=key, difficulty=6, diff_parts={K1 S3 I2 C0 F1}, diff_note={Cyclic functional equation x -> 1-1/x}]
If $f(x)+f\left(1-\dfrac1x\right)=1+x\ \forall\,x\in R-\{0,1\}$, then value of $f(2)$ is
\end{bankq}
\begin{bankq}[id=TOAS08-P2-MATH-S3-Q08, source={TOAS - Test 08 (15 May 2022)}, paper=2, qno=8, type=numerical, subtopic={Inverse}, answer={3.75}, answer_src=key, difficulty=7, diff_parts={K2 S3 I1 C1 F1}, diff_note={Inverse of the Mobius-linear composite function}]
Let $f$ be a real valued function such that $f\left(\dfrac{2x-3}{x-2}\right)=5x-2$, $x\neq2$, then value of $\dfrac54f^{-1}(13)$ is
\end{bankq}
\begin{bankpassage}[id=TOAS09-P1-MATH-S3-G3-4, source={TOAS - Test 09 (30 May 2022)}, questions={3-4}]
Let $f(x)=\sin^{-1}(\sin e^x)$, $g(x)=\sin^{-1}\!\left(\dfrac{3^x+3^{-x}}{4}\right)$ and $h(x)=f(x)-g(x)$
\end{bankpassage}
\begin{bankq}[id=TOAS09-P1-MATH-S3-Q03, source={TOAS - Test 09 (30 May 2022)}, paper=1, qno=3, type=numerical, subtopic={Range}, answer={8.00}, answer_src=key, difficulty=7, diff_parts={K2 S3 I1 C1 F1}, diff_note={Range: standard approach for this subtopic; 2 concept(s), 3 step(s), needs a non-obvious insight, heavier computation}, group=TOAS09-P1-MATH-S3-G3-4]
If range of $g(x)$ is $\left[\dfrac{\pi}{a},\dfrac{\pi}{b}\right]$, then value of $(a+b)$ is
\end{bankq}
\begin{bankq}[id=TOAS09-P1-MATH-S3-Q04, source={TOAS - Test 09 (30 May 2022)}, paper=1, qno=4, type=numerical, subtopic={Range}, answer={2.00}, answer_src=key, difficulty=7, diff_parts={K2 S3 I1 C1 F1}, diff_note={Range: standard approach for this subtopic; 2 concept(s), 3 step(s), needs a non-obvious insight, heavier computation}, group=TOAS09-P1-MATH-S3-G3-4]
Number of solution(s) of $h(x)=0$ is
\end{bankq}
\begin{bankq}[id=TOAS09-P2-MATH-S1-Q03, source={TOAS - Test 09 (30 May 2022)}, paper=2, qno=3, type=single, subtopic={Composition}, answer={A}, answer_src=key, difficulty=1, diff_parts={K1 S1 I0 C0 F0}, diff_note={Composition: standard approach for this subtopic; 1 concept(s), 1 step(s)}]
Which of the following is true-
\begin{qoptions}
\opt{A}{If number of elements in set A are finite, such that $f:A\to A$ is an one-one function, then $f$ is necessarily onto.}
\opt{B}{If a function is continuous in its domain and changes its sign for any 2 values of $x$ then odd number of roots lie between the given values of $x$}
\opt{C}{If $f:A\to A$ is one-one then it must be onto}
\opt{D}{It is possible that a curve contains a local maxima and a global minima at the same point.}
\end{qoptions}
\end{bankq}
\begin{bankq}[id=TOAS09-P2-MATH-S1-Q11, source={TOAS - Test 09 (30 May 2022)}, paper=2, qno=11, type=single, subtopic={Domain}, answer={D}, answer_src=key, difficulty=1, diff_parts={K1 S1 I0 C0 F0}, diff_note={Domain: standard approach for this subtopic; 1 concept(s), 1 step(s)}]
If $f(x)=3-x^2$, $1\leq x\leq4$, then the domain of $\log_e(f(2x))$ is
\begin{qoptions}
\opt{A}{$\left(\dfrac{-\sqrt3}{2},\dfrac{\sqrt3}{2}\right)$}
\opt{B}{$\left[\dfrac12,2\right]$}
\opt{C}{$\left(0,\dfrac{\sqrt3}{2}\right)$}
\opt{D}{$\left[\dfrac12,\dfrac{\sqrt3}{2}\right)$}
\end{qoptions}
\end{bankq}
\begin{bankq}[id=TOAS09-P2-MATH-S1-Q13, source={TOAS - Test 09 (30 May 2022)}, paper=2, qno=13, type=single, subtopic={Modulus and greatest integer}, answer={A}, answer_src=key, difficulty=4, diff_parts={K2 S2 I1 C0 F0}, diff_note={Modulus and greatest integer: standard approach for this subtopic; 2 concept(s), 2 step(s), needs a non-obvious insight}]
Number of positive integral values of `$k$' for which the equation $k=|x+|2x-1||-|x-|2x-1||$ has exactly three real solutions, is
\begin{qoptions}
\opt{A}{0}
\opt{B}{2}
\opt{C}{3}
\opt{D}{5}
\end{qoptions}
\end{bankq}
\begin{bankq}[id=TOAS09-P2-MATH-S2-Q09, source={TOAS - Test 09 (30 May 2022)}, paper=2, qno=9, type=numerical, subtopic={Range}, answer={33}, answer_src=solved, working={Denominator $\ge1$ so range $(0,100/3]$: integers 1 to 33}, difficulty=8, diff_parts={K3 S3 I1 C1 F1}, diff_note={Range: standard approach for this subtopic; 3 concept(s), 3 step(s), needs a non-obvious insight, heavier computation}]
Number of integers in the range of function $f(x)=\dfrac{100}{3\left(x^6+12x^4+5x^2+1\right)}$ is
\end{bankq}
\end{document}
